% XSEC: arXiv submission draft
% Compile: pdflatex xsec-submission.tex && pdflatex xsec-submission.tex

\documentclass[11pt]{article}

\usepackage[margin=1in]{geometry}
\usepackage[utf8]{inputenc}
\usepackage[T1]{fontenc}
\usepackage{hyperref}
\usepackage{url}
\usepackage{booktabs}
\usepackage{amsmath}
\usepackage{amsfonts}
\usepackage{microtype}
\usepackage{enumitem}
\usepackage{graphicx}
\usepackage{xcolor}
\usepackage[numbers]{natbib}

\hypersetup{
  colorlinks=true,
  linkcolor=blue!60!black,
  citecolor=blue!60!black,
  urlcolor=blue!60!black,
}

\title{XSEC: Methodology-Aware Evaluation and Cost-Aware Triage for Autonomous Pentesting Agents}

\author{
  XSEC Labs \\
  \texttt{github.com/uncesaii/xsec}
}

\date{}

\begin{document}

\maketitle

\begin{abstract}
Autonomous pentesting agents are frequently summarized by a single benchmark
headline, although reported performance is sensitive to retry protocol,
benchmark substrate, model/runtime, turn budget, and evidence policy. We
present XSEC, an open-source shell-first pentesting framework with layered
triage and blind verification, and pair it with a methodology-aware reporting
protocol that separates retained artifact-backed evidence from historical mixed
publication lines. In the current public ledger snapshot (2026-04-10), XSEC
reports a retained artifact-backed XBOW aggregate of 99/104. A 21-run triage
ablation (2026-04-11) shows mode-dependent behavior: full-moat triage is a
precision/recall tradeoff in white-box slices, a strict efficiency improvement
in black-box slices, and weakly additive on npm-bench over default scaffolding.
The primary contribution is a combined systems-and-methodology perspective in
which evidence lineage and protocol disclosure are treated as first-class
engineering surfaces for non-deterministic security agents.
\end{abstract}

\section{Introduction}

Single benchmark percentages for autonomous security agents hide variance and
protocol sensitivity. In practice, score changes can come from retry strategy,
benchmark fork, or evidence accounting policy rather than meaningful capability
changes. We frame XSEC as both:

\begin{itemize}[leftmargin=*]
  \item a \textbf{systems artifact}: shell-first exploitation, layered triage,
  blind verification, and multi-target operation;
  \item a \textbf{methodology artifact}: protocol-explicit benchmark reporting,
  retained-evidence lineage, and slice-aware interpretation.
\end{itemize}

\paragraph{Contributions.}
\begin{enumerate}[leftmargin=*,itemsep=2pt]
  \item A reproducible benchmarking protocol for non-deterministic pentesting
  agents (repeat framing, uncertainty, cost ceilings).
  \item A machine-auditable ledger model separating retained artifact-backed and
  historical mixed publication lines.
  \item A productionized multi-layer triage stack integrated in a live scanner
  pipeline, with per-layer verdict telemetry.
  \item A mode-dependent ablation narrative reporting both positive and negative
  outcomes, avoiding single-line ``moat works'' claims.
  \item A learned-routing foundation for per-finding policy selection,
  reported here as an offline modeling snapshot rather than a deployment
  claim: the 99/104 aggregate in Table~\ref{tab:ledger} is produced
  without the router active.
\end{enumerate}

\section{System Overview}

The pipeline is:
\[
\texttt{Plan} \rightarrow \texttt{Discover} \rightarrow \texttt{Attack} \rightarrow \texttt{Triage} \rightarrow \texttt{Verify} \rightarrow \texttt{Report}
\]

For web targets, XSEC uses shell-first execution (\texttt{bash} as primary
tool), with structured helpers optional. Findings then pass through layered
triage and verification before final reporting.

Core implementation paths include scanner orchestration, unified pipeline
execution, and triage modules in the core package.

\section{Benchmark Methodology}

We distinguish three reporting regimes:
\begin{itemize}
  \item single-shot,
  \item best-of-K,
  \item repeat-N per-attempt success rate with 95\% Wilson CI.
\end{itemize}

The benchmark ledger tracks:
\begin{itemize}
  \item retained artifact-backed totals (machine-reconstructible),
  \item historical mixed local+CI publication totals (legacy publication line).
\end{itemize}

This separation is implemented in
\texttt{packages/benchmark/results/benchmark-ledger.json} and reflected in the
docs benchmark surface.

\subsection{Current Ledger Snapshot (2026-04-10)}

\begin{table}[ht]
\centering
\small
\begin{tabular}{lccc}
\toprule
Surface & Black-box & White-box & Aggregate \\
\midrule
Retained artifact-backed & 74/104 & 79/104 & 99/104 \\
Historical mixed publication & 90/104 & 5 (white-box-only) & 95/104 \\
\bottomrule
\end{tabular}
\caption{Current ledger surfaces are intentionally separated and should not be merged.}
\label{tab:ledger}
\end{table}

\section{Triage Stack and Dynamic Routing}

The triage stack combines deterministic filters, feature extraction,
reachability and oracle checks, PoV/structured verify layers, memory-based
suppression, and optional debate/tree-search layers. Per-layer telemetry is
attached to findings, creating supervision data for learned per-finding routing.

Dynamic routing motivation is empirical: no static profile dominates across
white-box, black-box, and package slices.

\subsection{Ablation Snapshot (2026-04-11)}

\begin{table}[ht]
\centering
\small
\begin{tabular}{lrrrr}
\toprule
Profile & Flags & Findings & Cost (USD) & USD/flag \\
\midrule
\multicolumn{5}{l}{\textit{XBOW white-box (limit=50)}} \\
none & 43/50 & 67 & 14.34 & 0.33 \\
no-triage & 44/50 & 67 & 17.17 & 0.39 \\
moat-only & 41/50 & 25 & 26.89 & 0.66 \\
moat & 41/50 & 25 & 21.82 & 0.53 \\
\midrule
\multicolumn{5}{l}{\textit{XBOW black-box (limit=25)}} \\
none & 18/25 & 27 & 13.72 & 0.76 \\
no-triage & 19/25 & 34 & 10.37 & 0.55 \\
moat-only & 18/25 & 13 & 11.22 & 0.62 \\
moat & 19/25 & 14 & 10.04 & 0.53 \\
\bottomrule
\end{tabular}
\caption{Ablation snapshot showing mode-dependent moat behavior.}
\label{tab:ablation-xbow}
\end{table}

\begin{table}[ht]
\centering
\small
\begin{tabular}{lrrrr}
\toprule
Profile & F1 & TPR & FPR & Safe correct \\
\midrule
none & 0.973 & 1.00 & 0.11 & 24/27 \\
no-triage & 0.964 & 1.00 & 0.15 & 23/27 \\
moat-only & 0.964 & 1.00 & 0.15 & 23/27 \\
moat & 0.956 & 1.00 & 0.19 & 22/27 \\
default & 0.956 & 1.00 & 0.19 & 22/27 \\
\bottomrule
\end{tabular}
\caption{npm-bench profile outcomes (81 packages) from the same ablation set.}
\label{tab:ablation-npm}
\end{table}

\section{Empirical Highlights}

Current repository-reported posture includes:
\begin{itemize}
  \item retained artifact-backed XBOW aggregate at 99/104,
  \item separate black-box and white-box lines,
  \item a 21-run triage ablation matrix with mode-dependent outcomes.
\end{itemize}

Key methodological outcome: no static profile dominates all slices, so policy
selection should be context-aware.

\subsection{Dataset and router artifacts}

Current intermediate artifacts include:

\begin{itemize}
  \item triage dataset v1 stats: 969 samples,
  \item router v2 metadata: 1514 samples, 55 features, 5-fold CV mean
  F1\,$\approx$\,0.962.
\end{itemize}

These are treated as offline modeling artifacts, not final deployment claims.
The 99/104 retained artifact-backed aggregate reported in Table~\ref{tab:ledger}
is produced by the static-profile triage stack described above, with the
router v2 inactive at runtime.

\section{Evaluation Details and Interpretation}

This draft intentionally reports per-slice outcomes rather than collapsing all
settings into one global number.

\subsection{White-box versus black-box behavior}

The ablation in Table~\ref{tab:ablation-xbow} shows that the same triage
profile can have different utility depending on evidence availability. In
white-box mode, source access increases the attack agent's direct exploit signal
and some expensive triage layers can over-prune, producing a precision/recall
tradeoff. In black-box mode, where candidate findings are noisier, additional
verification layers provide clearer net benefit on findings efficiency and
USD/flag.

\subsection{Package slice caveat}

Table~\ref{tab:ablation-npm} shows behavior distinct from XBOW slices. The
current interpretation is that profile effects are entangled with scan
scaffolding and coarse package-level labels. This motivates treating package
evaluation as a separate operating regime rather than extrapolating web-target
behavior.

\subsection{Layer-level failure mode}

Single-layer isolation runs reported in the project docs identify one recurring
negative outlier (`egats`) on hard slices, while other layers are neutral or
positive in that context. This supports a policy of feature-gated layer rollout
and per-layer telemetry before broad default enablement.

\section{Discussion}

\subsection{Why methodology is part of the contribution}

For stochastic, tool-using security agents, benchmarking is not just
post-processing. Protocol choices alter reported results materially, so
measurement design and evidence accounting are engineering surfaces that deserve
explicit implementation and review.

\subsection{Operational policy implication}

The measured outcome is not "enable all layers always." A practical policy is:

\begin{enumerate}[leftmargin=*]
  \item keep deterministic/cheap layers broadly enabled,
  \item gate expensive layers by context,
  \item track per-layer impact over time,
  \item promote only slice-robust defaults.
\end{enumerate}

\subsection{From static profiles to routing}

The router path in current artifacts should be understood as an optimization
framework rather than a completed deployment claim. The offline metrics are
promising, but production value should be established by online A/B outcomes on
explicitly frozen benchmark protocols.

\section{Related Work}

Current XSEC docs position this work adjacent to three threads:

\begin{itemize}
  \item autonomous pentesting systems with differing orchestration styles
  (single-agent, role-based, and multi-agent workflows),
  \item triage pipelines combining deterministic checks and model-based
  judgments,
  \item hybrid feature+model signals for vulnerability classification and
  filtering.
\end{itemize}

The comparative caveat is central: cross-project benchmark numbers are only
loosely comparable unless protocol and substrate are aligned.

Representative systems in the current XSEC docs include BoxPwnr, Shannon,
KinoSec, Cyber-AutoAgent, deadend-cli, and
MAPTA\citep{boxpwnr,shannon,kinosec,cyberautoagent,deadend,mapta}.
For FP-triage architecture context, XSEC docs reference Endor Labs and
Semgrep Assistant disclosures\citep{endor,semgrep}. For hybrid feature+model
signals, XSEC docs use VulnBERT-style framing as design inspiration rather
than a direct task-equivalent claim\citep{vulnbert}.

Table~\ref{tab:competitors} collects the self-reported XBOW aggregates for
these systems. The numbers come from each project's own public disclosures,
not from a common evaluation run on matched hardware, model, or harness; they
are included for rough positioning, not as a leaderboard claim.

\begin{table}[ht]
\centering
\small
\begin{tabular}{llll}
\toprule
System & Aggregate & Model / Approach & Notes \\
\midrule
BoxPwnr\citep{boxpwnr}             & 97.1\% (101/104) & Best-of-N, 10 configs  & Context compaction, loop detection \\
Shannon\citep{shannon}              & 96.2\% (100/104) & White-box, multi-agent & Source-to-sink taint \\
\textbf{XSEC (this work)}         & \textbf{95.2\% (99/104)} & Shell-first, gpt-5.4   & Retained artifact-backed; see Table~\ref{tab:ledger} \\
KinoSec\citep{kinosec}              & 92.3\% (96/104)  & Black-box HTTP only    & 50-turn hard cap \\
Cyber-AutoAgent\citep{cyberautoagent} & 84.6\% (88/104)  & Meta-agent             & Self-rewriting prompts \\
deadend-cli\citep{deadend}          & 77.6\% ($\sim$76/98) & Single-agent CLI     & ADaPT decomposition \\
MAPTA\citep{mapta}                  & 76.9\% (80/104)  & 3-role multi-agent     & Evidence-gated branching \\
\bottomrule
\end{tabular}
\caption{Self-reported XBOW aggregates across representative autonomous pentesting systems. Cross-project numbers are protocol-sensitive (differing model, harness, retry policy, challenge subset) and should be interpreted as rough positioning rather than a matched-conditions leaderboard.}
\label{tab:competitors}
\end{table}

\section{Limitations and Threats to Validity}

We highlight benchmark dependence, run-to-run variance, slice sensitivity,
documentation and implementation drift risk, and cost-latency tradeoffs in deeper
triage pipelines.

Additional caveats:
\begin{itemize}
  \item package-level supervision can be coarse for per-finding labels;
  \item leaderboard comparisons across projects are protocol-sensitive and
  should be interpreted with caution;
  \item rapidly evolving repos can create temporary doc/code mismatch.
\end{itemize}

\section{Claim Traceability Checklist}

Before submission freeze, each numerical claim should be tied to:

\begin{enumerate}[leftmargin=*]
  \item an as-of date,
  \item an artifact path,
  \item protocol context (retry regime, budget, model/runtime),
  \item whether the claim is retained artifact-backed or historical mixed.
\end{enumerate}

The companion files \texttt{docs/paper/evaluation.md} and
\texttt{docs/paper/related\_work.md} exist to keep this audit step lightweight.

\section{Reproducibility Notes}

Canonical evidence files and runners are in:
\begin{itemize}
  \item \texttt{docs/src/content/docs/methodology.md}
  \item \texttt{docs/src/content/docs/benchmark.md}
  \item \texttt{docs/src/content/docs/research/2026-04-11-ablation.md}
  \item \texttt{docs/src/content/docs/research/fp-reduction-moat.md}
  \item \texttt{packages/benchmark/results/benchmark-ledger.json}
  \item \texttt{packages/benchmark/results/triage-dataset-v1.stats.json}
  \item \texttt{packages/benchmark/results/triage-router-v2-meta.json}
  \item \texttt{packages/benchmark/src/xbow-runner.ts}
  \item \texttt{packages/benchmark/src/wilson.ts}
\end{itemize}

Companion drafting notes:

\begin{itemize}
  \item \texttt{docs/paper/xsec.md}
  \item \texttt{docs/paper/evaluation.md}
  \item \texttt{docs/paper/related\_work.md}
\end{itemize}

\section{Conclusion}

For autonomous pentesting agents, measurement protocol quality is part of the
core contribution. XSEC's current evidence suggests that transparent protocol
disclosure, retained-evidence lineage, and mode-aware triage policy are more
reliable than single headline percentages.

\bibliographystyle{plainnat}
\bibliography{refs}

\end{document}
